\documentclass[10pt,a4paper]{amsart}
\usepackage[margin=19mm]{geometry}
\usepackage{amsmath,amssymb,mathtools}
\usepackage[scr=rsfs,cal=euler]{mathalfa}
\usepackage{xfrac}
\usepackage[unicode,hidelinks,bookmarks=false]{hyperref}
\newcommand{\ie}{i.e.\ }
\newcommand{\cf}{cf.\ }
\newcommand{\resp}{respectively\ }
\newcommand{\Subsection}{\subsection}
\providecommand{\nobreakdash}{\nobreak\hbox{-}}
\setlength{\emergencystretch}{2em}
\multlinegap0pt
\usepackage{xcolor}
\usepackage{amssymb}
\usepackage{bbm}
\usepackage{mathtools}
\usepackage{tikz}
\usepackage{enumerate}
\usepackage{breqn}
\newtheorem{lemma}{Lemma}
\title{On the existence of a morphism between certain Artin-Schreier curves}
\author{Beatriz Barbero Lucas, Stefano Lia, Gary McGuire}
\date{Source: arXiv:2602.15717v1 (CC BY 4.0)}
\begin{document}
\maketitle

\noindent\emph{Source and licence.} On the existence of a morphism between certain Artin-Schreier curves. Version arXiv:2602.15717v1 (CC BY 4.0), distributed by arXiv under the Creative Commons Attribution 4.0 International licence, http://creativecommons.org/licenses/by/4.0/, as recorded in the arXiv licence metadata for this exact version and shown on the abstract page https://arxiv.org/abs/2602.15717v1. Full licence notice: \texttt{licenses/arxiv-2602.15717-CC-BY-4.0.txt}.\par\smallskip

\noindent\textit{Editorial scope.} Counted unit: the article's lemma general\_div (source0.tex lines 890--909). The article's complete original proof of it is delivered with it (source0.tex lines 911--1002, one \texttt{proof} environment of 92 lines). No mathematics is written, repaired or re-derived: the statement and the proof are the article's own lines, in the article's own notation, and no retained line is substituted or re-flowed.  No further source-local context is needed: every cross-reference of every retained line resolves inside the two delivered spans.  Nothing was omitted from any retained span, so the package carries no omission record.  No retained line contains a citation, so the wrapper carries no bibliography and no bibliography entry.  No retained line contains a figure environment, an image-inclusion command or drawing code, so no image is delivered and the package declares no asset dependency.  Editorial formatting only, in addition: an AMS \texttt{amsart} preamble that restates the article's own declarations and macros used by the retained lines, namely the environments \texttt{lemma} on separate counters, together with the standard packages those lines need.  The retained environments print in this document as Lemma 1 in order of appearance, whereas the article prints the same items with the numbering of its own full document.  The complete native source of the article as distributed in the arXiv e-print for this exact version is bundled unchanged as \texttt{sources/194817/source0.tex}.
\begin{lemma}\label{general_div}
Let $p$ be an odd prime. Assume
\begin{enumerate}
  \item $k,\ell,r,t,m\in\mathbb{N}$, 
  \item $1<\ell < k$,  $0<r<k$, and $\ell < \frac{k-r}{2}$,
  \item $\ell\mid k$ and $k/\ell$ is even,
  \item $t\mid (p-1)$,
  \item $m\mid (p^k+1)(p-1)$ and $m>1$,
  \item $m\bigl(p^\ell t+1\bigr) \;=\; t\bigl(p^{k-r}+1\bigr).$
\end{enumerate}
Let $u=m/t$. Then 
\begin{enumerate}
\item $u\in \mathbb{Z}$ and $u\ge 5$.
\item If we write $u = 2^e u_0$ with $u_0$ odd,
then $u_0$ is a common divisor of $p^{k-r}+1$ and $p^k+1$.
\item $u_0$ divides $p^r-1$.
\item $\gcd(u_0,t)=1$.
\item $u \ge p^{(k-r)/2}$
\end{enumerate}
\end{lemma}

\begin{proof}
(1) \quad Note that
$  \gcd(p^\ell t+1,t) = 1$.
Therefore, from
\[
  m(p^\ell t+1) = t(p^{k-r}+1)
\]
we deduce that $t$ divides $m$. Let $u=m/t$.
Then
\begin{equation}\label{eq:uk}
  p^{k-r}+1 = u\,\bigl(p^\ell t+1\bigr)
\end{equation}
and it is clear that $u$ divides $p^{k-r}+1$.
In particular, $u_0$ divides $p^{k-r}+1$.

Next we are going to prove by contradiction that $u\neq 1,2,3,4$.

Suppose $u=1$, then \eqref{eq:uk} yields
$  p^{k-r}+1 = p^\ell t+1$,
so $p^{k-r} = p^\ell t$ and hence $t = p^{k-r-\ell}$. Since $k>r+\ell$,
we have $t\ge p$, which contradicts $t$ dividing $p-1$.
Thus $u\neq1$.

Let us suppose now that $u=2$, then from \eqref{eq:uk} we have
$p^{k-r}+1 = 2(p^\ell t+1)$,
and reducing this modulo $p$ gives $1\equiv 2 \pmod{p}$ which is a contradiction.
 Thus $u\neq2$. 

Suppose $u=3$, then from \eqref{eq:uk} we have
$ p^{k-r} = 3p^\ell t+2$,
and reducing this modulo $p$ gives $2\equiv 0 \pmod{p}$ which is a contradiction.
This proves that $u>3$.


Suppose $u=4$, then from \eqref{eq:uk} we have
$ p^{k-r} = 4p^\ell t+3$,
and reducing this modulo $p$ gives $3\equiv 0 \pmod{p}$ which is a contradiction if $p>3$.
This proves that $u>4$ when $p>3$.
If $u=4$ and $p=3$ the equation \eqref{eq:uk} becomes
$3^{k-r}+1 = 4\,\bigl(3^\ell t+1\bigr)$
or
$3^{k-r} = 4\cdot 3^\ell \cdot t+3$.
This is a contradiction because $k-r>1$ and $\ell >1$,
so mod 9 we get $0\equiv 3 \pmod{9}$.

(2) \quad  Write $u = 2^e u_0$ with $u_0$ odd.
It remains to prove that $u_0$ is a divisor of  $p^k+1$.

Write $p-1 = ts$ with $s\in\mathbb{N}$. From $m=ut$ and the
assumption $m\mid (p^k+1)(p-1)$ we have
\[
  ut  \mid (p^k+1)(p-1) = (p^k+1)ts,
\]
hence
\begin{equation}\label{eq:u-div}
  u \mid (p^k+1)s.
\end{equation}

Let $p_0$ be the odd part of $p^k+1$, and let $s_0$ be the odd part of $s$.
Then \eqref{eq:u-div} implies $u_0 | p_0 s_0$.
It is a known fact that $\gcd(p^k+1,p-1)=2$, which implies $\gcd(p_0,s_0)=1$.


If $p_1$ is the odd part of $p^{k-r}+1$, then 
\eqref{eq:uk} implies that $u_0$ divides $p_1$.
Also, $\gcd(p_1,s_0)=1$, by the fact that $\gcd(p^{k-r}+1,p-1)=2$.

If $\gcd(u_0,s_0)>1$ then this common divisor would divide 
both $s_0$ and $p_1$, a contradiction.
Therefore $\gcd(u_0,s_0)=1$, and $u_0 | p_0 s_0$ implies $u_0$ divides $p_0$.

As $u_0|p_0$ and $u_0|p_1$, we have obtained that $u_0$ is a common divisor of $p^{k-r}+1$ and $p^k+1$.  

(3) \quad  Since $u_0$ divides $p^k+1$ and $p^{k-r}+1$,
$u_0$ divides $p^r(p^{k-r}+1)-(p^k+1)=p^r-1$.

(4) \quad Recall that $t$ divides $p-1$ and $u_0$ divides $p_0$.
Then
\[
\gcd(p^k+1,p-1)=2 
\implies \gcd(p_0,p-1)=1 
\implies \gcd(u_0,t)=1.
\]


(5) \quad Since $t<p$ we get $tp^\ell < p^{\ell +1} \le p^{(k-r)/2}$
using the assumption $\ell < \frac{k-r}{2}$.

From the equation $p^{k-r}+1 = u\,\bigl(p^\ell t+1\bigr)$
we then get $p^{k-r}+1 \le u \cdot p^{(k-r)/2}$.
It follows that $u \ge p^{(k-r)/2}$.
\end{proof}

\end{document}
